\chapter*{Abstract} %l'asterisco dopo chapter serve per visualizzare il capitolo come "non numerato"

I Digital Twin sono una tecnologia emergente e trasversale, il loro impiego spazia dalla manutenzione predittiva di macchine industriali, alla gestione di flotte di agenti mobili autonomi. In molti contesti applicativi, la gestione e coordinazione di molteplici e differenti agenti presenta una sfida non banale. Questa tesi propone la implementazione di un framework per la realizzazione di complessi sistemi multi-agente di Digital Twin e distribuita, fondato sullo stato dell'arte presentato da carissimi colleghi e professori nella pubblicazione \textit{Advancing Robotic Systems with Distributed Multi-Agent Digital Twins: A Scalable and Adaptive Framework} \cite{Russo2025}.



% La crescente disponibilità nel mercato odierno di soluzioni robotiche già pronte, \textit{off-the-shelf},
%  destinate al pubblico e sempre più performanti,
% ha reso accessibile lo sviluppo e la condivisione di progetti relativi alla robotica 
% come mai prima d'ora.
% 
% La robotica ripercorre un percorso già visto nell'ambito dello sviluppo di videogiochi,
% con il primo framework nel 2005 per il game development con licenza gratuita.
% 
% Il presente lavoro accoglie le ultime novità nel game development e nella robotica accessibile,
% attraverso la progettazione di un videogioco che impiega la tecnologia del Digital Twin, con l'obiettivo
% di esplorare nuove forme di interazione tra mondo reale e virtuale, promuovendo al contempo un'esperienza divertente e innovativa.
% A tale scopo, mediante il potente ecosistema di sviluppo offerto da \textit{Espressif} per l'Alvik,
% la libreria \textit{enet} per le comunicazioni, e il versatile motore di gioco \textit{Godot}, 
% è stato prodotto un videogioco multi giocatore, dove quattro player possono sfidarsi
% all'interno di una \textit{arena virtuale e fisica}.
% 
% Nel videogioco, ogni player comanda un carro, tale carro è intepretato nel mondo reale da un
% Alvik, il quale condividerà con il suo gemello virtuale informazioni sull'ambiente circostante. \\
% Ogni carro è dotato di cannone laser e tre palloncini, l'obiettivo è quello di far scoppiare i
% palloncini degli avversari: vince l'ultimo giocatore in gioco dotato di palloncini.
% 